\chapter{Connective tissue diseases}
\index{Connective tissue diseases}

\medskip
\noindent\textit{\textbf{Under review.} This chapter is an AI-assisted educational draft; it carries no source citations and is not a substitute for professional medical advice.}

\section*{Definition and Overview}
Connective tissue diseases are a group of disorders in which the immune system attacks the body's connective tissue, the material that supports and binds organs, joints, skin, and blood vessels. The main types are autoimmune and are sometimes called systemic autoimmune diseases or collagen vascular diseases. The most common are rheumatoid arthritis, systemic lupus erythematosus (SLE), Sjögren's syndrome, scleroderma, and the inflammatory muscle diseases. They can affect many parts of the body at once, and symptoms often come and go over many years.

\section*{Epidemiology}
Connective tissue diseases as a group are not rare. Rheumatoid arthritis affects roughly 1 in 100 people, and systemic lupus erythematosus affects roughly 1 in 1000, with a strong female predominance (about 9 in 10 cases). Most begin between the ages of 20 and 50, though some types appear in children or older adults. Scleroderma and the inflammatory myopathies are less common. Certain ethnic groups have higher rates of lupus, and a family history of autoimmune disease increases the risk.

\section*{Aetiology and Pathophysiology}
The exact cause is unknown, but these diseases develop when the immune system mistakenly produces antibodies against the body's own connective tissue:
\begin{itemize}
    \item \textbf{Genetics:} certain genes increase susceptibility, though most affected people have no affected relatives
    \item \textbf{Hormonal factors:} the strong female predominance suggests a hormonal role
    \item \textbf{Environmental triggers:} viral infections, sunlight, smoking, and some medicines may trigger disease in susceptible people
    \item \textbf{Immune dysregulation:} autoantibodies and inflammation damage joints, skin, blood vessels, and internal organs
    \item \textbf{Inflammation:} chronic inflammation leads to tissue damage, scarring (fibrosis), and impaired organ function
\end{itemize}

\section*{Clinical Features}
Symptoms vary by disease but commonly include:
\begin{itemize}
    \item Joint pain, swelling, and stiffness, especially in the morning (rheumatoid arthritis)
    \item Fatigue, fever, and feeling generally unwell
    \item Skin rashes, including the malar (butterfly) rash across the cheeks in lupus, and hardened, tight skin in scleroderma
    \item Sensitivity to sunlight with rashes or flares
    \item Dry eyes and dry mouth (Sjögren's syndrome)
    \item Raynaud's phenomenon: fingers turning white, blue, then red in the cold
    \item Hair loss, mouth ulcers, and swollen glands
    \item Muscle weakness and pain in the inflammatory myopathies
\end{itemize}

\section*{Differential Diagnosis}
\begin{itemize}
    \item Osteoarthritis, which is wear-and-tear rather than inflammatory
    \item Fibromyalgia, causing widespread pain without visible inflammation
    \item Infection, including viral arthritis or chronic infections
    \item Chronic fatigue syndrome
    \item Endocrine disorders such as thyroid disease
    \item Vasculitis (inflammation of blood vessels alone)
    \item Cancer, which can rarely mimic systemic disease
\end{itemize}

\section*{When to Seek Medical Care}
See a doctor if you have persistent joint swelling, an unexplained rash, profound fatigue, or symptoms affecting several body systems together.

\textbf{Contact your local emergency services for:}
\begin{itemize}
    \item Chest pain or difficulty breathing
    \item Sudden weakness, confusion, or a fit
    \item Severe headache with vision changes
    \item Sudden cold, painful, or numb fingers or toes in a finger that turns blue or black
\end{itemize}

\section*{Diagnosis and Investigations}
\begin{itemize}
    \item \textbf{History and examination:} a detailed account of symptoms across joints, skin, and organs, looking for patterns characteristic of each disease
    \item \textbf{Blood tests:} inflammatory markers (ESR, CRP), full blood count, and kidney and liver function
    \item \textbf{Autoantibody tests:} rheumatoid factor, anti-CCP, antinuclear antibody (ANA), anti-dsDNA, and others help identify the specific disease
    \item \textbf{Urine testing:} to detect kidney involvement in lupus
    \item \textbf{Imaging:} X-rays, ultrasound, or MRI of affected joints; chest CT for lung involvement
    \item \textbf{Biopsy:} skin, muscle, or kidney biopsy to confirm the diagnosis
    \item \textbf{Specialist referral:} rheumatologist-led assessment, with other specialists as organs are involved
\end{itemize}

\section*{Management}
\begin{itemize}
    \item \textbf{Medicines:} non-steroidal anti-inflammatory drugs (NSAIDs) and corticosteroids for rapid symptom control; disease-modifying drugs such as methotrexate; biologic and targeted therapies for moderate to severe disease
    \item \textbf{Symptom control:} analgesics for pain, medicines for fatigue, and treatments for dry eyes and mouth
    \item \textbf{Sun protection:} high-SPF sunscreen and protective clothing for light-sensitive disease
    \item \textbf{Allied health:} physiotherapy and occupational therapy to maintain joint and muscle function
    \item \textbf{Lifestyle:} balanced rest and activity, a healthy diet, and smoking cessation
    \item \textbf{Pregnancy planning:} specialist review, as some treatments are not safe in pregnancy
    \item \textbf{Regular monitoring:} blood tests and organ checks to detect complications early
\end{itemize}

\section*{Complications}
\begin{itemize}
    \item Joint damage and deformity from long-standing inflammation
    \item Kidney disease (nephritis), especially in lupus
    \item Heart and blood vessel disease, including accelerated atherosclerosis and inflammation of the heart lining
    \item Lung involvement, including scarring and pulmonary hypertension
    \item Anaemia and reduced immunity from the disease or its treatment
    \item Osteoporosis from long-term corticosteroid use
    \item Increased risk of infection with immunosuppressive treatment
\end{itemize}

\section*{Prognosis and Outcomes}
Connective tissue diseases are usually chronic and require long-term care, but the outlook has improved greatly with modern treatment. Many people with rheumatoid arthritis or lupus lead full lives with good symptom control, though some have progressive organ damage. Disease patterns vary widely: some people have mild, intermittent symptoms, while others have severe, active disease. Early diagnosis, specialist care, and adherence to treatment reduce complications. Regular review remains important because the disease and its treatment both change over time.

\section*{Prevention and Self-Care}
\begin{itemize}
    \item Stop smoking, which worsens many connective tissue diseases
    \item Protect skin from the sun with high-SPF sunscreen and protective clothing
    \item Stay physically active within the limits set by symptoms
    \item Learn about the disease and build a good relationship with your rheumatology team
    \item Take medicines as prescribed and report side effects rather than stopping them
    \item Keep up to date with recommended vaccinations (as advised, avoiding live vaccines when immunosuppressed)
    \item Manage stress, and seek support for fatigue, mood, and day-to-day challenges
    \item Attend regular reviews and blood tests to catch complications early
\end{itemize}
